    \expectedproblemsection{Direct Calculation Problems}

    \begin{expectedproblemgroup}{Finite Sums and Products}
        Evaluate each of the following finite sums or products exactly. State any
        restrictions on the integer parameters that are needed for the displayed
        expression to be defined.

        \begin{expectedsubproblem}[Vandermonde's Identity]
            For $n\in\Nzero$, determine
            \[
                \sum_{k=0}^{n}\binom{n}{k}^2.
            \]
        \end{expectedsubproblem}

        \begin{expectedsubsolution}
            This is the $r=s=n$ specialization of Vandermonde's identity from
            Mathematical Concepts. Its coefficient interpretation gives the
            same result immediately: since
            \[
                \binom{n}{k}=\binom{n}{n-k},
            \]
            the required sum is the coefficient of $x^n$ in
            \[
                (1+x)^n(1+x)^n=(1+x)^{2n}.
            \]
            Therefore,
            \[
                \boxed{
                    \sum_{k=0}^{n}\binom{n}{k}^2=\binom{2n}{n}
                }.
            \]
        \end{expectedsubsolution}

        \begin{expectedsubproblem}
            For $n\in\Nzero$, determine
            \[
                \sum_{k=0}^{n}
                (-1)^k\binom{n}{k}\frac{1}{k+1}.
            \]
        \end{expectedsubproblem}

        \begin{expectedsubsolution}
            Using
            \[
                \frac{1}{k+1}=\int_0^1x^k\dd x,
            \]
            we obtain
            \[
                \begin{aligned}
                    \sum_{k=0}^{n}
                    (-1)^k\binom{n}{k}\frac{1}{k+1}
                    &=
                    \int_0^1
                    \sum_{k=0}^{n}\binom{n}{k}(-x)^k\dd x\\
                    &=
                    \int_0^1(1-x)^n\dd x\\
                    &=
                    \boxed{\frac{1}{n+1}}.
                \end{aligned}
            \]
        \end{expectedsubsolution}

        \begin{expectedsubproblem}[Sine Product Formula]
            For an integer $n\geq2$, determine
            \[
                \prod_{k=1}^{n-1}\sin\left(\frac{k\pi}{n}\right).
            \]
        \end{expectedsubproblem}

        \begin{expectedsubsolution}
            Let $\omega_n=e^{2\pi i/n}$. Factoring $z^n-1$ and then dividing by
            $z-1$ gives
            \[
                \frac{z^n-1}{z-1}
                =
                \prod_{k=1}^{n-1}(z-\omega_n^k).
            \]
            Letting $z\to1$ and taking absolute values yields
            \[
                n
                =
                \prod_{k=1}^{n-1}|1-\omega_n^k|
                =
                \prod_{k=1}^{n-1}
                2\sin\left(\frac{k\pi}{n}\right).
            \]
            Hence
            \[
                \boxed{
                    \prod_{k=1}^{n-1}
                    \sin\left(\frac{k\pi}{n}\right)
                    =
                    \frac{n}{2^{n-1}}
                }.
            \]
        \end{expectedsubsolution}

        \begin{expectedsubproblem}
            For an integer $n\geq2$, determine
            \[
                \prod_{k=2}^{n}\frac{k^3-1}{k^3+1}.
            \]
        \end{expectedsubproblem}

        \begin{expectedsubsolution}
            Factor each term as
            \[
                \frac{k^3-1}{k^3+1}
                =
                \frac{k-1}{k+1}
                \frac{k^2+k+1}{k^2-k+1}.
            \]
            The two products telescope separately:
            \[
                \prod_{k=2}^{n}\frac{k-1}{k+1}
                =
                \frac{2}{n(n+1)},
            \]
            and, since $k^2+k+1=(k+1)^2-(k+1)+1$,
            \[
                \prod_{k=2}^{n}
                \frac{k^2+k+1}{k^2-k+1}
                =
                \frac{n^2+n+1}{3}.
            \]
            Therefore,
            \[
                \boxed{
                    \prod_{k=2}^{n}\frac{k^3-1}{k^3+1}
                    =
                    \frac{2(n^2+n+1)}{3n(n+1)}
                }.
            \]
        \end{expectedsubsolution}

        \begin{expectedsubproblem}
            For $m\in\N$ and $\sin x\neq0$, determine
            \[
                \prod_{k=0}^{m-1}\cos(2^k x).
            \]
        \end{expectedsubproblem}

        \begin{expectedsubsolution}
            Repeatedly apply
            \[
                \sin(2t)=2\sin t\cos t.
            \]
            This gives
            \[
                \sin(2^m x)
                =
                2^m\sin x
                \prod_{k=0}^{m-1}\cos(2^k x).
            \]
            Since $\sin x\neq0$,
            \[
                \boxed{
                    \prod_{k=0}^{m-1}\cos(2^k x)
                    =
                    \frac{\sin(2^m x)}{2^m\sin x}
                }.
            \]
        \end{expectedsubsolution}

        \begin{expectedsubproblem}
            For $n\in\Nzero$, determine
            \[
                \sum_{\substack{0\leq k\leq n\\3\mid k}}
                \binom{n}{k}.
            \]
        \end{expectedsubproblem}

        \begin{expectedsubsolution}
            Let
            \[
                \omega=e^{2\pi i/3}.
            \]
            Since
            \[
                1+\omega^k+\omega^{2k}
                =
                \begin{cases}
                    3, & 3\mid k,\\
                    0, & 3\nmid k,
                \end{cases}
            \]
            the roots-of-unity filter gives
            \[
                \begin{aligned}
                    3S_n
                    &=
                    \sum_{k=0}^{n}\binom{n}{k}
                    \bigl(1+\omega^k+\omega^{2k}\bigr)\\
                    &=
                    2^n+(1+\omega)^n+(1+\omega^2)^n.
                \end{aligned}
            \]
            Now
            \[
                1+\omega=e^{i\pi/3},
                \qquad
                1+\omega^2=e^{-i\pi/3}.
            \]
            Therefore,
            \[
                \boxed{
                    S_n
                    =
                    \frac{2^n+2\cos(n\pi/3)}{3}
                }.
            \]
        \end{expectedsubsolution}

        \begin{expectedsubproblem}[1997 CSAT]
            For a permutation
            \[
                (a_1,a_2,a_3,a_4)
            \]
            of $\{1,2,3,4\}$, let $s_k$ be the number of entries to the right
            of $a_k$ that are smaller than $a_k$, for $k=1,2,3$. Define
            \[
                |(a_1,a_2,a_3,a_4)|=s_1+s_2+s_3.
            \]
            Determine the sum of this quantity over all $24$ permutations.
        \end{expectedsubproblem}

        \begin{expectedsubsolution}
            For each unordered pair of distinct elements, exactly half of the
            $24$ permutations place the larger element before the smaller one.
            Thus each of the
            \[
                \binom42=6
            \]
            pairs contributes one inversion in exactly $12$ permutations. Hence
            the total inversion number is
            \[
                \boxed{6\cdot12=72}.
            \]
        \end{expectedsubsolution}

    \end{expectedproblemgroup}

    \begin{expectedproblemgroup}{Recurrence Relations}
        Determine the requested term, closed form, or residue for each of the
        following recursively defined sequences.

        \begin{expectedsubproblem}[2006 VTRMC]
            Let $F_0=0$, $F_1=1$, and
            \[
                F_{n+2}=F_{n+1}+F_n.
            \]
            Determine the last decimal digit of $F_{2006}$.
        \end{expectedsubproblem}

        \begin{expectedsubsolution}
            Modulo $2$, the Fibonacci sequence has period $3$:
            \[
                0,1,1,0,1,1,\ldots.
            \]
            Hence
            \[
                F_{2006}\equiv F_2\equiv1\pmod2.
            \]
            Also
            \[
                F_{n+5}=5F_{n+1}+3F_n,
            \]
            so modulo $5$,
            \[
                F_{n+5}\equiv3F_n.
            \]
            Iterating four times gives $F_{n+20}\equiv F_n\pmod5$. Therefore
            \[
                F_{2006}\equiv F_6=8\equiv3\pmod5.
            \]
            The unique decimal digit that is odd and congruent to $3$ modulo $5$ is
            \[
                \boxed{3}.
            \]
        \end{expectedsubsolution}


        \begin{expectedsubproblem}
            Let $a_0=1$ and
            \[
                a_{n+1}
                =
                \frac12\left(a_n+\frac{2}{a_n}\right).
            \]
            Determine a closed form for $a_n$.
        \end{expectedsubproblem}

        \begin{expectedsubsolution}
            Set
            \[
                r_n=\frac{a_n-\sqrt2}{a_n+\sqrt2}.
            \]
            Direct substitution into the recurrence gives
            \[
                r_{n+1}=r_n^2.
            \]
            Moreover,
            \[
                r_0
                =
                \frac{1-\sqrt2}{1+\sqrt2}
                =
                -(3-2\sqrt2).
            \]
            Thus
            \[
                r_n=\bigl(-(3-2\sqrt2)\bigr)^{2^n}.
            \]
            Solving the defining relation for $a_n$ yields
            \[
                \boxed{
                    a_n
                    =
                    \sqrt2\,
                    \frac{1+r_n}{1-r_n},
                    \qquad
                    r_n=\bigl(-(3-2\sqrt2)\bigr)^{2^n}
                }.
            \]
        \end{expectedsubsolution}


        \begin{expectedsubproblem}
            Let $a_0=0$ and
            \[
                a_{n+1}=\frac{1+a_n}{2+a_n}.
            \]
            Determine a closed form for $a_n$ in terms of the Fibonacci
            sequence.
        \end{expectedsubproblem}

        \begin{expectedsubsolution}
            Let $(F_n)$ be the Fibonacci sequence with $F_0=0$ and $F_1=1$. We
            claim that
            \[
                a_n=\frac{F_{2n}}{F_{2n+1}}.
            \]
            The formula holds for $n=0$. If it holds for $n$, then
            \[
                \begin{aligned}
                    a_{n+1}
                    &=
                    \frac{F_{2n}+F_{2n+1}}
                    {F_{2n}+2F_{2n+1}}\\
                    &=
                    \frac{F_{2n+2}}{F_{2n+3}}.
                \end{aligned}
            \]
            Thus
            \[
                \boxed{a_n=\frac{F_{2n}}{F_{2n+1}}}.
            \]
        \end{expectedsubsolution}

        \begin{expectedsubproblem}[2014 \UTokyo{} Entrance Examination]
            A bag initially contains $a+2$ white balls and one red ball, where
            $a$ is a positive integer. After each draw, perform the following
            operation:
            \begin{problemchoices}
                \item if a white ball is drawn, replace the contents by $a$
                white balls and one red ball;
                \item if the red ball is drawn, remove it and leave the other
                balls unchanged.
            \end{problemchoices}
            Let $p_n$ be the probability that the red ball is drawn on the
            $n$th draw. Determine $p_n$.
        \end{expectedsubproblem}

        \begin{expectedsubsolution}
            The first probability is
            \[
                p_1=\frac{1}{a+3}.
            \]
            If the red ball is drawn on the $n$th draw, it cannot be drawn next. If
            a white ball is drawn, the next state has $a+1$ balls, exactly one of
            which is red. Therefore
            \[
                p_{n+1}=\frac{1-p_n}{a+1}.
            \]
            The fixed point of this recurrence is $1/(a+2)$. Hence, with
            \[
                q_n=p_n-\frac{1}{a+2},
            \]
            we have
            \[
                q_{n+1}=-\frac{q_n}{a+1},
                \qquad
                q_1=-\frac{1}{(a+3)(a+2)}.
            \]
            Consequently,
            \[
                \boxed{
                    p_n
                    =
                    \frac{1}{a+2}
                    +
                    \frac{(-1)^n}{
                        (a+3)(a+2)(a+1)^{n-1}
                    }
                }
            \]
            for $n\geq1$. In particular,
            \[
                p_2=\frac{a+2}{(a+3)(a+1)}.
            \]
        \end{expectedsubsolution}

    \end{expectedproblemgroup}

    \begin{expectedproblemgroup}{Limits}
        Evaluate each of the following limits. In parameterized problems,
        determine the exact value for every parameter regime stated in the
        subproblem.

        \begin{expectedsubproblem}[2012 \KAIST{} Problem of the Week]
            Let $a_0=3$ and
            \[
                a_n=a_{n-1}+\sqrt{a_{n-1}^2+3},
                \qquad n\geq1.
            \]
            Determine
            \[
                \lim\limits_{n\to\infty}\frac{a_n}{2^n}.
            \]
        \end{expectedsubproblem}

        \begin{expectedsubsolution}
            Write
            \[
                a_n=\sqrt3\cot\theta_n.
            \]
            Since $a_0=3$, we may take $\theta_0=\pi/6$. The half-angle identity
            \[
                \cot\frac{\theta}{2}=\cot\theta+\csc\theta
            \]
            transforms the recurrence into
            \[
                a_n
                =
                \sqrt3\left(\cot\theta_{n-1}+\csc\theta_{n-1}\right)
                =
                \sqrt3\cot\frac{\theta_{n-1}}{2}.
            \]
            Hence
            \[
                \theta_n=\frac{\pi}{6\cdot2^n}
                \qquad\text{and}\qquad
                a_n=\sqrt3\cot\left(\frac{\pi}{6\cdot2^n}\right).
            \]
            Using $t\cot t\to1$ as $t\to0$, we obtain
            \[
                \boxed{
                    \lim\limits_{n\to\infty}\frac{a_n}{2^n}
                    =
                    \frac{6\sqrt3}{\pi}
                }.
            \]
        \end{expectedsubsolution}

        \begin{expectedsubproblem}[2013 \KAIST{} Problem of the Week]
            For $a,b\in\R$, determine
            \[
                \lim\limits_{n\to\infty}
                n\left(
                    1-\frac{a}{n}-\frac{b\ln(n+1)}{n}
                \right)^n.
            \]
        \end{expectedsubproblem}

        \begin{expectedsubsolution}
            Put
            \[
                c_n=a+b\ln(n+1).
            \]
            Since $c_n/n\to0$ and $c_n^2/n\to0$, the logarithmic expansion gives
            \[
                n\ln\left(1-\frac{c_n}{n}\right)
                =
                -c_n+o(1).
            \]
            Therefore,
            \[
                n\left(1-\frac{c_n}{n}\right)^n
                =
                e^{-a}n(n+1)^{-b}\bigl(1+o(1)\bigr).
            \]
            It follows that the limit is
            \[
                \boxed{
                    \begin{cases}
                        \infty, & b<1,\\
                        e^{-a}, & b=1,\\
                        0, & b>1.
                    \end{cases}
                }
            \]
        \end{expectedsubsolution}

        \begin{expectedsubproblem}[2014 \KAIST{} Problem of the Week]
            For $x\geq0$, let
            \[
                f_n(x)
                =
                \frac{
                    \displaystyle\prod_{k=1}^{n-1}(x+k)(x+k+1)
                }{(n!)^2}.
            \]
            Determine
            \[
                \lim\limits_{n\to\infty}f_n(x).
            \]
        \end{expectedsubproblem}

        \begin{expectedsubsolution}
            The Gamma function gives
            \[
                f_n(x)
                =
                \frac{\Gamma(n+x)\Gamma(n+x+1)}
                {\Gamma(x+1)\Gamma(x+2)\Gamma(n+1)^2}.
            \]
            Using
            \[
                \frac{\Gamma(n+\alpha)}{\Gamma(n+\beta)}
                \sim n^{\alpha-\beta},
            \]
            we obtain
            \[
                f_n(x)
                \sim
                \frac{n^{2x-1}}{\Gamma(x+1)\Gamma(x+2)}.
            \]
            Consequently,
            \[
                \boxed{
                    \lim\limits_{n\to\infty}f_n(x)
                    =
                    \begin{cases}
                        0, & 0\leq x<\dfrac12,\\[4pt]
                        \dfrac{8}{3\pi}, & x=\dfrac12,\\[6pt]
                        \infty, & x>\dfrac12.
                    \end{cases}
                }
            \]
        \end{expectedsubsolution}

        \begin{expectedsubproblem}[2016 \KAIST{} Problem of the Week]
            For $a\geq0$, determine
            \[
                \lim\limits_{n\to\infty}
                n\int_{-1}^{0}
                \left(x+\frac{x^2}{2}+e^{ax}\right)^n\dd x.
            \]
        \end{expectedsubproblem}

        \begin{expectedsubsolution}
            Let
            \[
                g(x)=x+\frac{x^2}{2}+e^{ax}.
            \]
            On $[-1,0]$, we have $-1/2\leq g(x)\leq1$, with equality $g(x)=1$
            only at $x=0$. Moreover,
            \[
                g(0)=1
                \qquad\text{and}\qquad
                g'(0)=1+a.
            \]
            The contribution from every interval $[-1,-\delta]$ is exponentially
            small. On the remaining interval, set $x=-t/n$. For each fixed $t\geq0$,
            \[
                g\left(-\frac{t}{n}\right)^n
                \longrightarrow
                e^{-(1+a)t}.
            \]
            A uniform exponential bound near the endpoint permits dominated
            convergence, and therefore
            \[
                \begin{aligned}
                    \lim\limits_{n\to\infty}
                    n\int_{-1}^{0}g(x)^n\dd x
                    &=
                    \int_0^\infty e^{-(1+a)t}\dd t\\
                    &=
                    \boxed{\frac{1}{1+a}}.
                \end{aligned}
            \]
        \end{expectedsubsolution}

        \begin{expectedsubproblem}[2023 \KAIST{} Problem of the Week]
            Determine
            \[
                \lim\limits_{n\to\infty}
                \left(
                    \frac{\displaystyle\sum_{k=1}^{n+2}k^k}
                    {\displaystyle\sum_{k=1}^{n+1}k^k}
                    -
                    \frac{\displaystyle\sum_{k=1}^{n+1}k^k}
                    {\displaystyle\sum_{k=1}^{n}k^k}
                \right).
            \]
        \end{expectedsubproblem}

        \begin{expectedsubsolution}
            Let
            \[
                S_n=\sum_{k=1}^{n}k^k.
            \]
            The final two terms dominate the sum:
            \[
                S_n
                =
                n^n\left(1+\frac{1}{en}+O\left(\frac{1}{n^2}\right)\right).
            \]
            Hence
            \[
                \begin{aligned}
                    \frac{S_{n+1}}{S_n}
                    &=
                    (n+1)\left(1+\frac{1}{n}\right)^n
                    \frac{
                        1+\dfrac{1}{e(n+1)}+O(n^{-2})
                    }{
                        1+\dfrac{1}{en}+O(n^{-2})
                    }\\
                    &=
                    en+\frac{e}{2}+O\left(\frac{1}{n}\right).
                \end{aligned}
            \]
            Replacing $n$ by $n+1$ and subtracting gives
            \[
                \boxed{
                    \lim\limits_{n\to\infty}
                    \left(
                        \frac{S_{n+2}}{S_{n+1}}
                        -
                        \frac{S_{n+1}}{S_n}
                    \right)
                    =e
                }.
            \]
        \end{expectedsubsolution}

    \end{expectedproblemgroup}

    \begin{expectedproblemgroup}{Derivatives}
        Compute each requested derivative or derivative-dependent quantity
        exactly. For high-order derivatives, a closed form involving standard
        sequences or harmonic numbers is acceptable.

        \begin{expectedsubproblem}[2006 HMMT]
            A nonzero polynomial $f\in\R[x]$ satisfies
            \[
                f(x)=f'(x)f''(x).
            \]
            Determine the leading coefficient of $f$.
        \end{expectedsubproblem}

        \begin{expectedsubsolution}
            Suppose that the leading term of $f$ is $cx^n$, where $c\neq0$. The
            leading term of $f'f''$ is
            \[
                c^2n^2(n-1)x^{2n-3}.
            \]
            Comparing degrees gives
            \[
                n=2n-3,
            \]
            so $n=3$. Comparing leading coefficients then gives
            \[
                c=18c^2.
            \]
            Since $c\neq0$,
            \[
                \boxed{c=\frac{1}{18}}.
            \]
        \end{expectedsubsolution}




        \begin{expectedsubproblem}
            Let
            \[
                H_n=\sum_{k=1}^{n}\frac{1}{k}.
            \]
            Determine
            \[
                \left.
                \frac{\dd^{100}}{\dd x^{100}}
                \bigl(x^{100}\ln x\bigr)
                \right|_{x=1}.
            \]
        \end{expectedsubproblem}

        \begin{expectedsubsolution}
            For a real parameter $\alpha$,
            \[
                \frac{\dd^{100}}{\dd x^{100}}x^\alpha
                =
                \alpha(\alpha-1)\cdots(\alpha-99)x^{\alpha-100}.
            \]
            Differentiate both sides with respect to $\alpha$ and then set
            $\alpha=100$ and $x=1$. Since
            \[
                \left.
                \frac{\dd}{\dd\alpha}
                \prod_{j=0}^{99}(\alpha-j)
                \right|_{\alpha=100}
                =
                100!\sum_{k=1}^{100}\frac{1}{k},
            \]
            we obtain
            \[
                \boxed{100!\,H_{100}}.
            \]
        \end{expectedsubsolution}

    \end{expectedproblemgroup}

    \begin{expectedproblemgroup}{Series}
        Determine whether each of the following series converges. For every
        convergent series, determine its exact sum. In parameterized problems,
        state all parameter values for which convergence occurs.

        \begin{expectedsubproblem}[2016 \KAIST{} Problem of the Week]
            Let $k$ be a positive integer, and define
            \[
                a_n
                =
                \begin{cases}
                    1, & (k+1)\nmid n,\\
                    -k, & (k+1)\mid n.
                \end{cases}
            \]
            Determine
            \[
                \sum_{n=1}^{\infty}\frac{a_n}{n}.
            \]
        \end{expectedsubproblem}

        \begin{expectedsubsolution}
            At the end of the $N$th complete block of length $k+1$, the partial sum
            is
            \[
                \begin{aligned}
                    \sum_{n=1}^{(k+1)N}\frac{a_n}{n}
                    &=
                    \sum_{n=1}^{(k+1)N}\frac{1}{n}
                    -(k+1)\sum_{m=1}^{N}\frac{1}{(k+1)m}\\
                    &=
                    H_{(k+1)N}-H_N.
                \end{aligned}
            \]
            Since $H_N=\ln N+\gamma+o(1)$, these partial sums tend to
            $\ln(k+1)$. The terms inside the final incomplete block tend uniformly
            to zero, so the full series has the same limit. Therefore,
            \[
                \boxed{
                    \sum_{n=1}^{\infty}\frac{a_n}{n}=\ln(k+1)
                }.
            \]
        \end{expectedsubsolution}

        \begin{expectedsubproblem}[2017 \KAIST{} Problem of the Week]
            Determine whether
            \[
                \sum_{n=0}^{\infty}
                \frac{1}{4^n}\binom{2n}{n}
            \]
            converges.
        \end{expectedsubproblem}

        \begin{expectedsubsolution}
            By the central binomial coefficient asymptotic,
            \[
                \frac{1}{4^n}\binom{2n}{n}
                \sim
                \frac{1}{\sqrt{\pi n}}.
            \]
            Since the $p$-series $\sum n^{-1/2}$ diverges, the limit comparison
            test gives
            \[
                \boxed{
                    \sum_{n=0}^{\infty}
                    \frac{1}{4^n}\binom{2n}{n}
                    \text{ diverges}
                }.
            \]
        \end{expectedsubsolution}

        \begin{expectedsubproblem}[2014 \KAIST{} Problem of the Week]
            Determine all positive integers $\ell$ for which
            \[
                \sum_{n=1}^{\infty}
                \frac{n^3}{(n+1)(n+2)\cdots(n+\ell)}
            \]
            converges, and compute its value whenever it converges.
        \end{expectedsubproblem}

        \begin{expectedsubsolution}
            The $n$th term is asymptotic to $n^{3-\ell}$, so the series converges
            exactly when
            \[
                \ell\geq5.
            \]
            For such $\ell$, the beta-integral identity gives
            \[
                \frac{1}{(n+1)(n+2)\cdots(n+\ell)}
                =
                \frac{1}{(\ell-1)!}
                \int_0^1x^n(1-x)^{\ell-1}\dd x.
            \]
            Since
            \[
                \sum_{n=1}^{\infty}n^3x^n
                =
                \frac{x(1+4x+x^2)}{(1-x)^4},
            \]
            term-by-term integration yields
            \[
                \begin{aligned}
                    \sum_{n=1}^{\infty}
                    \frac{n^3}{(n+1)\cdots(n+\ell)}
                    &=
                    \frac{1}{(\ell-1)!}
                    \int_0^1x(1+4x+x^2)(1-x)^{\ell-5}\dd x\\
                    &=
                    \boxed{
                        \frac{\ell(\ell+5)}{
                            (\ell-4)(\ell-3)(\ell-2)(\ell-1)(\ell-1)!
                        }
                    }.
                \end{aligned}
            \]
        \end{expectedsubsolution}

        \begin{expectedsubproblem}
            Determine
            \[
                \sum_{n=1}^{\infty}
                \frac{1}{n^2\binom{2n}{n}}.
            \]
        \end{expectedsubproblem}

        \begin{expectedsubsolution}
            The Taylor expansion
            \[
                2\arcsin^2\left(\frac{x}{2}\right)
                =
                \sum_{n=1}^{\infty}
                \frac{x^{2n}}{n^2\binom{2n}{n}}
            \]
            holds for $|x|\leq2$. It may be verified by differentiating both sides
            and using the power series for $1/\sqrt{1-x^2/4}$.

            Setting $x=1$ gives
            \[
                \sum_{n=1}^{\infty}
                \frac{1}{n^2\binom{2n}{n}}
                =
                2\arcsin^2\left(\frac12\right)
                =
                2\left(\frac{\pi}{6}\right)^2
                =
                \boxed{\frac{\pi^2}{18}}.
            \]
        \end{expectedsubsolution}

        \begin{expectedsubproblem}
            Evaluate
            \[
                \sum_{n=1}^{\infty}\frac{n}{(n+1)!}
            \]
            in two ways.
        \end{expectedsubproblem}

        \begin{expectedsubsolution}
            First,
            \[
                \frac{n}{(n+1)!}
                =
                \frac{1}{n!}-\frac{1}{(n+1)!},
            \]
            so the series telescopes:
            \[
                \sum_{n=1}^{N}\frac{n}{(n+1)!}
                =
                1-\frac{1}{(N+1)!}
                \longrightarrow
                1.
            \]

            Alternatively, put $m=n+1$ and use the exponential series:
            \[
                \begin{aligned}
                    \sum_{n=1}^{\infty}\frac{n}{(n+1)!}
                    &=
                    \sum_{m=2}^{\infty}\frac{m-1}{m!}\\
                    &=
                    \sum_{m=2}^{\infty}\frac{1}{(m-1)!}
                    -
                    \sum_{m=2}^{\infty}\frac{1}{m!}\\
                    &=
                    (e-1)-(e-2)=1.
                \end{aligned}
            \]
            Thus both methods give
            \[
                \boxed{1}.
            \]
        \end{expectedsubsolution}

        \begin{expectedsubproblem}[1998 CSAT]
            Let
            \[
                a_1=1,
                \qquad
                a_2=2,
                \qquad
                a_{n+2}=a_{n+1}+a_n.
            \]
            Determine
            \[
                \sum_{n=1}^{\infty}
                \frac{a_n}{a_{n+1}a_{n+2}}.
            \]
        \end{expectedsubproblem}

        \begin{expectedsubsolution}
            The recurrence gives
            \[
                \frac{a_n}{a_{n+1}a_{n+2}}
                =
                \frac{a_{n+2}-a_{n+1}}{a_{n+1}a_{n+2}}
                =
                \frac{1}{a_{n+1}}-\frac{1}{a_{n+2}}.
            \]
            Thus the series telescopes:
            \[
                \sum_{n=1}^{N}
                \frac{a_n}{a_{n+1}a_{n+2}}
                =
                \frac{1}{a_2}-\frac{1}{a_{N+2}}.
            \]
            Since $a_n\to\infty$,
            \[
                \boxed{
                    \sum_{n=1}^{\infty}
                    \frac{a_n}{a_{n+1}a_{n+2}}
                    =
                    \frac12
                }.
            \]
        \end{expectedsubsolution}

    \end{expectedproblemgroup}

    \begin{expectedproblemgroup}{Integrals}
        Evaluate each of the following integrals. For indefinite integrals, the
        constant of integration may be omitted; for definite integrals,
        determine the exact value.

        \begin{expectedsubproblem}[2010 \MIT{} Integration Bee]
            \(\displaystyle
                \int_0^{\pi/2}\sin x\sin(2x)\sin(3x)\dd x
            \)
        \end{expectedsubproblem}

        \begin{expectedsubsolution}
            By the product-to-sum identities,
            \[
                \sin x\sin(3x)
                =
                \frac{1}{2}\bigl(\cos(2x)-\cos(4x)\bigr),
            \]
            and hence
            \[
                \sin x\sin(2x)\sin(3x)
                =
                \frac{1}{4}
                \bigl(\sin(2x)+\sin(4x)-\sin(6x)\bigr).
            \]
            Therefore,
            \[
                \int_0^{\pi/2}\sin x\sin(2x)\sin(3x)\dd x
                =
                \frac{1}{4}\left(1-\frac{1}{3}\right)
                =
                \boxed{\frac{1}{6}}.
            \]
        \end{expectedsubsolution}

        \begin{expectedsubproblem}[2010 \MIT{} Integration Bee]
            \(\displaystyle
                \int_0^1\sin^2(\ln x)\dd x
            \)
        \end{expectedsubproblem}

        \begin{expectedsubsolution}
            Let $t=-\ln x$. Then $x=e^{-t}$ and
            $\dd x=-e^{-t}\dd t$, so
            \[
                \int_0^1\sin^2(\ln x)\dd x
                =
                \int_0^\infty e^{-t}\sin^2t\dd t.
            \]
            Since $\sin^2t=(1-\cos(2t))/2$,
            \[
                \int_0^\infty e^{-t}\sin^2t\dd t
                =
                \frac{1}{2}
                \left(1-\frac{1}{5}\right)
                =
                \boxed{\frac{2}{5}}.
            \]
        \end{expectedsubsolution}

        \begin{expectedsubproblem}[2010 \MIT{} Integration Bee]
            \(\displaystyle
                \int_1^\infty\frac{\dd x}{x\sqrt{x^4-1}}
            \)
        \end{expectedsubproblem}

        \begin{expectedsubsolution}
            Let $u=x^2$. Then
            \[
                \frac{\dd x}{x}=\frac{\dd u}{2u},
            \]
            and therefore
            \[
                \int_1^\infty\frac{\dd x}{x\sqrt{x^4-1}}
                =
                \frac{1}{2}
                \int_1^\infty\frac{\dd u}{u\sqrt{u^2-1}}.
            \]
            With $u=\sec\theta$, this becomes
            \[
                \frac{1}{2}\int_0^{\pi/2}\dd\theta
                =
                \boxed{\frac{\pi}{4}}.
            \]
        \end{expectedsubsolution}

        \begin{expectedsubproblem}[2010 \MIT{} Integration Bee]
            \(\displaystyle
                \int_0^1\frac{\ln(1+x)}{1+x^2}\dd x
            \)
        \end{expectedsubproblem}

        \begin{expectedsubsolution}
            Let
            \[
                I=\int_0^1\frac{\ln(1+x)}{1+x^2}\dd x.
            \]
            Substituting $x=\tan t$ gives
            \[
                I=\int_0^{\pi/4}\ln(1+\tan t)\dd t.
            \]
            Under $t\mapsto\pi/4-t$,
            \[
                1+\tan\left(\frac{\pi}{4}-t\right)
                =
                \frac{2}{1+\tan t}.
            \]
            Thus
            \[
                I=\frac{\pi}{4}\ln2-I,
            \]
            and consequently
            \[
                \boxed{I=\frac{\pi}{8}\ln2}.
            \]
        \end{expectedsubsolution}

        \begin{expectedsubproblem}[2010 \MIT{} Integration Bee]
            \(\displaystyle
                \int\frac{\cos x+x\sin x}{x(x+\cos x)}\dd x
            \)
        \end{expectedsubproblem}

        \begin{expectedsubsolution}
            We recognize the logarithmic derivative
            \[
                \diff{}{x}
                \bigl(\ln|x|-\ln|x+\cos x|\bigr)
                =
                \frac{\cos x+x\sin x}{x(x+\cos x)}.
            \]
            Hence
            \[
                \boxed{
                    \int\frac{\cos x+x\sin x}{x(x+\cos x)}\dd x
                    =
                    \ln|x|-\ln|x+\cos x|
                }.
            \]
        \end{expectedsubsolution}

        \begin{expectedsubproblem}[2010 \MIT{} Integration Bee]
            \(\displaystyle
                \int_0^{\pi/2}\frac{\dd x}{\sin x+\sec x}
            \)
        \end{expectedsubproblem}

        \begin{expectedsubsolution}
            Write
            \[
                I
                =
                \int_0^{\pi/2}
                \frac{\cos x}{1+\sin x\cos x}\dd x.
            \]
            Replacing $x$ by $\pi/2-x$ also gives
            \[
                I
                =
                \int_0^{\pi/2}
                \frac{\sin x}{1+\sin x\cos x}\dd x.
            \]
            Adding the two expressions and setting
            $u=\sin x-\cos x$, we obtain
            \[
                2I
                =
                \int_{-1}^{1}\frac{2}{3-u^2}\dd u.
            \]
            Therefore,
            \[
                \boxed{
                    I=\frac{1}{\sqrt3}\ln(2+\sqrt3)
                }.
            \]
        \end{expectedsubsolution}

        \begin{expectedsubproblem}[2010 \MIT{} Integration Bee]
            \(\displaystyle
                \int_0^1\sqrt{1+x\sqrt{1+x\sqrt{1+\cdots}}}\dd x
            \)
        \end{expectedsubproblem}

        \begin{expectedsubsolution}
            For $x\in[0,1]$, define the finite truncations by
            \[
                R_1(x)=1,
                \qquad
                R_{n+1}(x)=\sqrt{1+xR_n(x)}.
            \]
            Since $y\mapsto\sqrt{1+xy}$ is increasing for $x\in[0,1]$, induction
            gives
            \[
                1\leq R_n(x)\leq\varphi
                \qquad\text{and}\qquad
                R_n(x)\leq R_{n+1}(x),
            \]
            where $\varphi=(1+\sqrt5)/2$ and $\varphi^2=1+\varphi$. Hence
            $R_n(x)$ converges pointwise to a limit $R(x)$. Passing to the limit in
            the recurrence gives
            \[
                R(x)^2=1+xR(x),
            \]
            so positivity yields
            \[
                R(x)=\frac{x+\sqrt{x^2+4}}{2}.
            \]
            The uniform bound $0\leq R_n(x)\leq\varphi$ permits integration by the
            dominated convergence theorem. Therefore,
            \[
                \begin{aligned}
                    \int_0^1R(x)\dd x
                    &=
                    \frac12\int_0^1x\dd x
                    +
                    \frac12\int_0^1\sqrt{x^2+4}\dd x\\
                    &=
                    \boxed{
                        \frac{1+\sqrt5}{4}
                        +
                        \ln\left(\frac{1+\sqrt5}{2}\right)
                    }.
                \end{aligned}
            \]
        \end{expectedsubsolution}

        \begin{expectedsubproblem}[2010 \MIT{} Integration Bee]
            \(\displaystyle
                \int\left(\frac{1}{\ln x}-\frac{1}{(\ln x)^2}\right)\dd x,
                \qquad x>0,\quad x\neq1
            \)
        \end{expectedsubproblem}

        \begin{expectedsubsolution}
            On any interval contained in $(0,\infty)\setminus\{1\}$,
            \[
                \diff{}{x}\left(\frac{x}{\ln x}\right)
                =
                \frac{1}{\ln x}
                -
                \frac{1}{(\ln x)^2}.
            \]
            Therefore,
            \[
                \boxed{
                    \int
                    \left(
                        \frac{1}{\ln x}
                        -
                        \frac{1}{(\ln x)^2}
                    \right)\dd x
                    =
                    \frac{x}{\ln x}
                }.
            \]
        \end{expectedsubsolution}

        \begin{expectedsubproblem}[2010 \MIT{} Integration Bee]
            \(\displaystyle
                \int_1^2\sqrt{x-1}\sqrt{2-x}\dd x
            \)
        \end{expectedsubproblem}

        \begin{expectedsubsolution}
            Let $u=x-1$. Then
            \[
                \int_1^2\sqrt{x-1}\sqrt{2-x}\dd x
                =
                \int_0^1\sqrt{u(1-u)}\dd u.
            \]
            Setting $v=u-1/2$ gives
            \[
                \sqrt{u(1-u)}
                =
                \sqrt{\frac{1}{4}-v^2}.
            \]
            Thus the integral is the area of a semicircle of radius $1/2$:
            \[
                \boxed{\frac{\pi}{8}}.
            \]
        \end{expectedsubsolution}

        \begin{expectedsubproblem}[2011 \MIT{} Integration Bee]
            \(\displaystyle
                \int\cos(\ln x)\dd x,
                \qquad x>0
            \)
        \end{expectedsubproblem}

        \begin{expectedsubsolution}
            Let $t=\ln x$. Then $x=e^t$ and $\dd x=e^t\dd t$, so
            \[
                \int\cos(\ln x)\dd x
                =
                \int e^t\cos t\dd t.
            \]
            Therefore,
            \[
                \boxed{
                    \int\cos(\ln x)\dd x
                    =
                    \frac{x}{2}
                    \bigl(\cos(\ln x)+\sin(\ln x)\bigr)
                }.
            \]
        \end{expectedsubsolution}

        \begin{expectedsubproblem}[2011 \MIT{} Integration Bee]
            \(\displaystyle
                \int_0^\pi\cos x\cos(3x)\cos(5x)\dd x
            \)
        \end{expectedsubproblem}

        \begin{expectedsubsolution}
            By repeated use of the product-to-sum identity,
            \[
                \cos x\cos(3x)\cos(5x)
                =
                \frac{1}{4}
                \bigl(
                    \cos x+\cos(3x)+\cos(7x)+\cos(9x)
                \bigr).
            \]
            Each term has integral $0$ on $[0,\pi]$. Hence
            \[
                \boxed{
                    \int_0^\pi\cos x\cos(3x)\cos(5x)\dd x=0
                }.
            \]
        \end{expectedsubsolution}

        \begin{expectedsubproblem}[2011 \MIT{} Integration Bee]
            \(\displaystyle
                \int\sin(101x)\sin^{99}x\dd x
            \)
        \end{expectedsubproblem}

        \begin{expectedsubsolution}
            Differentiating a suitable product gives
            \[
                \begin{aligned}
                    \diff{}{x}
                    \left(
                        \frac{1}{100}\sin(100x)\sin^{100}x
                    \right)
                    &=
                    \sin^{99}x
                    \bigl(
                        \sin x\cos(100x)
                        +
                        \cos x\sin(100x)
                    \bigr)\\
                    &=
                    \sin(101x)\sin^{99}x.
                \end{aligned}
            \]
            Therefore,
            \[
                \boxed{
                    \int\sin(101x)\sin^{99}x\dd x
                    =
                    \frac{1}{100}\sin(100x)\sin^{100}x
                }.
            \]
        \end{expectedsubsolution}

        \begin{expectedsubproblem}[2011 \MIT{} Integration Bee]
            \(\displaystyle
                \int x e^{e^{x^2}+x^2}\dd x
            \)
        \end{expectedsubproblem}

        \begin{expectedsubsolution}
            Since
            \[
                e^{e^{x^2}+x^2}=e^{e^{x^2}}e^{x^2},
            \]
            let $u=e^{x^2}$. Then $\dd u=2xe^{x^2}\dd x$, and
            \[
                \int x e^{e^{x^2}+x^2}\dd x
                =
                \frac{1}{2}\int e^u\dd u.
            \]
            Hence
            \[
                \boxed{
                    \int x e^{e^{x^2}+x^2}\dd x
                    =
                    \frac{1}{2}e^{e^{x^2}}
                }.
            \]
        \end{expectedsubsolution}

        \begin{expectedsubproblem}[2011 \MIT{} Integration Bee]
            \(\displaystyle
                \int\sqrt{\frac{1-x}{1+x}}\dd x,
                \qquad -1<x<1
            \)
        \end{expectedsubproblem}

        \begin{expectedsubsolution}
            For $-1<x<1$,
            \[
                \sqrt{\frac{1-x}{1+x}}
                =
                \frac{1-x}{\sqrt{1-x^2}}.
            \]
            Therefore,
            \[
                \begin{aligned}
                    \int\sqrt{\frac{1-x}{1+x}}\dd x
                    &=
                    \int\frac{\dd x}{\sqrt{1-x^2}}
                    -
                    \int\frac{x}{\sqrt{1-x^2}}\dd x\\
                    &=
                    \boxed{\arcsin x+\sqrt{1-x^2}}.
                \end{aligned}
            \]
        \end{expectedsubsolution}

        \begin{expectedsubproblem}[Gabriel's Horn]
            Determine the volume and the area of the curved surface obtained by
            rotating
            \[
                y=\frac{1}{x},
                \qquad
                x\geq1,
            \]
            about the $x$-axis.
        \end{expectedsubproblem}

        \begin{expectedsubsolution}
            The volume is
            \[
                V
                =
                \pi\int_1^{\infty}\frac{1}{x^2}\dd x
                =
                \boxed{\pi}.
            \]
            The curved surface area is
            \[
                S
                =
                2\pi\int_1^{\infty}
                \frac{1}{x}\sqrt{1+\frac{1}{x^4}}\dd x.
            \]
            Since the integrand is at least $1/x$, this integral dominates
            \[
                2\pi\int_1^{\infty}\frac{1}{x}\dd x,
            \]
            which diverges. Hence
            \[
                \boxed{S=\infty}.
            \]
        \end{expectedsubsolution}

        \begin{expectedsubproblem}[Napkin Ring Problem]
            A cylindrical hole is drilled through the center of a sphere. The
            remaining solid has height $6$. Determine its volume without being
            given either the sphere radius or the hole radius.
        \end{expectedsubproblem}

        \begin{expectedsubsolution}
            Let the original sphere radius be $R$ and the cylindrical hole radius
            be $r$. If the remaining solid has height $6$, then its half-height is
            $3$, so
            \[
                R^2-r^2=9.
            \]
            At height $z\in[-3,3]$, the cross-section of the remaining solid is an
            annulus of area
            \[
                \pi\bigl((R^2-z^2)-r^2\bigr)
                =
                \pi(9-z^2).
            \]
            Therefore,
            \[
                V
                =
                \pi\int_{-3}^{3}(9-z^2)\dd z
                =
                \boxed{36\pi}.
            \]
        \end{expectedsubsolution}

        \begin{expectedsubproblem}[Steinmetz Solid]
            Two infinite circular cylinders of radius $R$ have perpendicular
            axes that intersect. Determine the volume of their intersection.
        \end{expectedsubproblem}

        \begin{expectedsubsolution}
            Choose coordinates so that the cylinders are
            \[
                x^2+y^2\leq R^2,
                \qquad
                x^2+z^2\leq R^2.
            \]
            For fixed $x\in[-R,R]$, both $y$ and $z$ range independently over
            \[
                \left[-\sqrt{R^2-x^2},\sqrt{R^2-x^2}\right].
            \]
            The cross-section is therefore a square of area
            \[
                4(R^2-x^2).
            \]
            Hence
            \[
                V
                =
                \int_{-R}^{R}4(R^2-x^2)\dd x
                =
                \boxed{\frac{16}{3}R^3}.
            \]
        \end{expectedsubsolution}

        \begin{expectedsubproblem}[Dirichlet Integral]
            Determine the improper integral
            \[
                \int_0^{\infty}\frac{\sin x}{x}\dd x.
            \]
        \end{expectedsubproblem}

        \begin{expectedsubsolution}
            For $a>0$, define
            \[
                F(a)
                =
                \int_0^{\infty}e^{-ax}\frac{\sin x}{x}\dd x.
            \]
            Differentiation under the integral sign gives
            \[
                F'(a)
                =
                -\int_0^{\infty}e^{-ax}\sin x\dd x
                =
                -\frac{1}{1+a^2}.
            \]
            Also $F(a)\to0$ as $a\to\infty$. Therefore,
            \[
                F(a)
                =
                \int_a^{\infty}\frac{\dd t}{1+t^2}
                =
                \frac{\pi}{2}-\arctan a.
            \]
            Letting $a\to0^+$ and applying Abel's limiting argument for the
            Dirichlet integral yields
            \[
                \boxed{
                    \int_0^{\infty}\frac{\sin x}{x}\dd x
                    =
                    \frac{\pi}{2}
                }.
            \]
        \end{expectedsubsolution}

        \begin{expectedsubproblem}[1944 Dalzell, On 22/7]
            Evaluate
            \[
                \int_0^1
                \frac{x^4(1-x)^4}{1+x^2}\dd x,
            \]
            and use the result to prove that
            \[
                \pi<\frac{22}{7}.
            \]
        \end{expectedsubproblem}

        \begin{expectedsubsolution}
            Polynomial division gives
            \[
                \frac{x^4(1-x)^4}{1+x^2}
                =
                x^6-4x^5+5x^4-4x^2+4
                -
                \frac{4}{1+x^2}.
            \]
            Therefore,
            \[
                \begin{aligned}
                    I
                    &=
                    \int_0^1
                    \left(x^6-4x^5+5x^4-4x^2+4\right)\dd x
                    -4\int_0^1\frac{\dd x}{1+x^2}\\
                    &=
                    \frac{22}{7}-\pi.
                \end{aligned}
            \]
            Since the original integrand is positive on $(0,1)$,
            \[
                0<I=\frac{22}{7}-\pi,
            \]
            and hence
            \[
                \boxed{\pi<\frac{22}{7}}.
            \]
        \end{expectedsubsolution}

        \begin{expectedsubproblem}[Borwein Integrals]
            For $n\in\Nzero$, define
            \[
                I_n
                =
                \int_0^{\infty}
                \prod_{k=0}^{n}
                \frac{\sin(x/(2k+1))}{x/(2k+1)}\dd x.
            \]
            Determine how long the pattern $I_n=\pi/2$ persists, and show that
            the next case is smaller.
        \end{expectedsubproblem}

        \begin{expectedsubsolution}
            Put
            \[
                a_k=\frac{1}{2k+1},
            \]
            and let $X_k$ be independent random variables uniformly distributed on
            $[-a_k,a_k]$. The characteristic function of $X_k$ is
            \[
                \frac{\sin(a_k t)}{a_k t}.
            \]
            If $S_n=X_0+\cdots+X_n$ has density $f_n$, Fourier inversion at the
            origin gives
            \[
                I_n=\pi f_n(0).
            \]

            Write $S_n=X_0+Y_n$, where $X_0$ is uniform on $[-1,1]$. If
            \[
                \sum_{k=1}^{n}a_k\leq1,
            \]
            then $|Y_n|\leq1$ always, and conditioning on $Y_n$ gives
            \[
                f_n(0)=\frac12.
            \]
            Thus $I_n=\pi/2$. Now
            \[
                \sum_{k=1}^{6}\frac{1}{2k+1}
                =
                \frac{43024}{45045}<1,
            \]
            whereas
            \[
                \sum_{k=1}^{7}\frac{1}{2k+1}
                =
                \frac{46027}{45045}>1.
            \]
            Hence
            \[
                \boxed{I_0=I_1=\cdots=I_6=\frac{\pi}{2}}.
            \]
            For $n=7$, the continuous random variable $Y_7$ has positive
            probability of satisfying $|Y_7|>1$. Therefore
            \[
                f_7(0)
                =
                \frac12\Prob(|Y_7|\leq1)
                <
                \frac12,
            \]
            and consequently
            \[
                \boxed{I_7<\frac{\pi}{2}}.
            \]
        \end{expectedsubsolution}

        \begin{expectedsubproblem}[1987 Putnam]
            Determine
            \[
                \int_2^4
                \frac{\ln(9-x)}{
                    \ln(9-x)+\ln(x+3)
                }
                \dd x.
            \]
        \end{expectedsubproblem}

        \begin{expectedsubsolution}
            Let
            \[
                h(x)
                =
                \frac{\ln(9-x)}{
                    \ln(9-x)+\ln(x+3)
                }.
            \]
            The substitution $x\mapsto6-x$ gives
            \[
                h(6-x)=1-h(x).
            \]
            Since the interval $[2,4]$ is symmetric about $3$,
            \[
                2\int_2^4h(x)\dd x
                =
                \int_2^4\bigl(h(x)+h(6-x)\bigr)\dd x
                =2.
            \]
            Therefore the integral is
            \[
                \boxed{1}.
            \]
        \end{expectedsubsolution}

        \begin{expectedsubproblem}[1989 \UTokyo{} Entrance Examination]
            Let
            \[
                f(x)=\pi x^2\sin(\pi x^2).
            \]
            Determine the volume obtained by revolving the region bounded by
            $y=f(x)$, the $x$-axis, $x=0$, and $x=1$ about the $y$-axis.
        \end{expectedsubproblem}

        \begin{expectedsubsolution}
            Using cylindrical shells, the volume is
            \[
                V
                =
                2\pi\int_0^1x f(x)\dd x
                =
                2\pi^2\int_0^1x^3\sin(\pi x^2)\dd x.
            \]
            Let $u=x^2$. Then
            \[
                V
                =
                \pi^2\int_0^1u\sin(\pi u)\dd u.
            \]
            Integration by parts gives
            \[
                \int_0^1u\sin(\pi u)\dd u=\frac1\pi.
            \]
            Hence
            \[
                \boxed{V=\pi}.
            \]
        \end{expectedsubsolution}

        \begin{expectedsubproblem}[2014 CSAT]
            Let $f\colon\R\to\R$ be continuous and odd. Suppose that
            \[
                f(x)
                =
                \frac{\pi}{2}
                \int_1^{x+1}f(t)\dd t
            \]
            for every real $x$, and that $f(1)=1$. Determine
            \[
                \pi^2\int_0^1x f(x+1)\dd x.
            \]
        \end{expectedsubproblem}

        \begin{expectedsubsolution}
            Differentiating the integral equation gives
            \[
                f'(x)=\frac{\pi}{2}f(x+1),
            \]
            so
            \[
                f(x+1)=\frac{2}{\pi}f'(x).
            \]
            Therefore the required expression is
            \[
                2\pi\int_0^1x f'(x)\dd x.
            \]
            Integration by parts yields
            \[
                2\pi\left(
                    [xf(x)]_0^1-\int_0^1f(x)\dd x
                \right)
                =
                2\pi-2\pi\int_0^1f(x)\dd x.
            \]

            Since $f$ is odd and $f(1)=1$, we have $f(-1)=-1$. Substituting
            $x=-1$ into the original equation gives
            \[
                -1
                =
                \frac{\pi}{2}\int_1^0f(t)\dd t
                =
                -\frac{\pi}{2}\int_0^1f(t)\dd t.
            \]
            Thus
            \[
                \int_0^1f(t)\dd t=\frac{2}{\pi},
            \]
            and the required value is
            \[
                \boxed{2\pi-4}.
            \]
        \end{expectedsubsolution}

        \begin{expectedsubproblem}
            For a real parameter $p>1$, evaluate
            \[
                \int_0^\infty\frac{\dd x}{1+x^p}.
            \]
        \end{expectedsubproblem}

        \begin{expectedsubsolution}
            Set
            \[
                t=\frac{x^p}{1+x^p}.
            \]
            Then
            \[
                x^p=\frac{t}{1-t}
            \]
            and the integral becomes
            \[
                \frac{1}{p}
                \int_0^1
                t^{1/p-1}(1-t)^{-1/p}\dd t
                =
                \frac{1}{p}
                B\left(\frac{1}{p},1-\frac{1}{p}\right).
            \]
            By the beta--gamma identity and Euler's reflection formula,
            \[
                \Gamma(z)\Gamma(1-z)=\frac{\pi}{\sin(\pi z)},
            \]
            so
            \[
                \boxed{
                    \int_0^\infty\frac{\dd x}{1+x^p}
                    =
                    \frac{\pi}{p}
                    \csc\left(\frac{\pi}{p}\right)
                }.
            \]
        \end{expectedsubsolution}

    \end{expectedproblemgroup}

    \begin{expectedproblemgroup}{Complex Numbers}
        Evaluate each of the following expressions exactly. When roots of unity
        occur, use $\omega_n=e^{2\pi i/n}$ unless another root is specified.



        \begin{expectedsubproblem}
            For an integer $n\geq2$, determine
            \[
                \sum_{k=0}^{n-1}(1+\omega_n^k)^n.
            \]
        \end{expectedsubproblem}

        \begin{expectedsubsolution}
            Expand by the binomial theorem and interchange the finite sums:
            \[
                \sum_{k=0}^{n-1}(1+\omega_n^k)^n
                =
                \sum_{j=0}^{n}\binom{n}{j}
                \sum_{k=0}^{n-1}\omega_n^{kj}.
            \]
            The inner sum is $n$ when $n\mid j$ and $0$ otherwise. For
            $0\leq j\leq n$, only $j=0$ and $j=n$ contribute. Hence
            \[
                \boxed{
                    \sum_{k=0}^{n-1}(1+\omega_n^k)^n=2n
                }.
            \]
        \end{expectedsubsolution}

        \begin{expectedsubproblem}[Quadratic Gauss Sum]
            Let $\omega=e^{2\pi i/5}$. Determine
            \[
                \sum_{k=0}^{4}\omega^{k^2}.
            \]
        \end{expectedsubproblem}

        \begin{expectedsubsolution}
            The quadratic residues modulo $5$ are
            \[
                0,1,4,4,1.
            \]
            Therefore,
            \[
                \sum_{k=0}^{4}\omega^{k^2}
                =
                1+2(\omega+\omega^{-1})
                =
                1+4\cos\left(\frac{2\pi}{5}\right).
            \]
            Since
            \[
                2\cos\left(\frac{2\pi}{5}\right)
                =
                \frac{\sqrt5-1}{2},
            \]
            the sum is
            \[
                \boxed{\sqrt5}.
            \]
        \end{expectedsubsolution}

        \begin{expectedsubproblem}
            For an integer $n\geq2$, determine
            \[
                \sum_{k=0}^{n-1}\frac{1}{2-\omega_n^k}.
            \]
        \end{expectedsubproblem}

        \begin{expectedsubsolution}
            Let
            \[
                p(z)=z^n-1
                =
                \prod_{k=0}^{n-1}(z-\omega_n^k).
            \]
            Its logarithmic derivative is
            \[
                \frac{p'(z)}{p(z)}
                =
                \sum_{k=0}^{n-1}\frac{1}{z-\omega_n^k}.
            \]
            Setting $z=2$ gives
            \[
                \boxed{
                    \sum_{k=0}^{n-1}\frac{1}{2-\omega_n^k}
                    =
                    \frac{n2^{n-1}}{2^n-1}
                }.
            \]
        \end{expectedsubsolution}

    \end{expectedproblemgroup}

    \begin{expectedproblemgroup}{Polynomials and Roots}
        Use the coefficients, factorization, or symmetric functions of the roots
        to compute each requested quantity exactly. Solving every polynomial
        explicitly is not required unless the problem asks for the roots.

        \begin{expectedsubproblem}[2006 HMMT]
            Determine the sum of the distinct real roots of
            \[
                2x^6-3x^5+3x^4+x^3-3x^2+3x-1=0.
            \]
        \end{expectedsubproblem}

        \begin{expectedsubsolution}
            The polynomial factors as
            \[
                2x^6-3x^5+3x^4+x^3-3x^2+3x-1
                =
                (x^2-x+1)^2(2x^2+x-1).
            \]
            The first factor has no real roots. The distinct real roots are therefore
            the roots of
            \[
                2x^2+x-1=(2x-1)(x+1),
            \]
            namely $1/2$ and $-1$. Their sum is
            \[
                \boxed{-\frac12}.
            \]
        \end{expectedsubsolution}

        \begin{expectedsubproblem}[1997 VTRMC]
            Two roots $r_1,r_2$ of
            \[
                x^4-x^3+ax^2-8x-8=0
            \]
            satisfy $r_1r_2=2$. Determine $a$, $r_1$, and $r_2$.
        \end{expectedsubproblem}

        \begin{expectedsubsolution}
            Since $r_1r_2=2$, their quadratic factor has the form
            \[
                x^2+px+2.
            \]
            The remaining factor must have constant term $-4$, so write
            \[
                x^4-x^3+ax^2-8x-8
                =
                (x^2+px+2)(x^2+qx-4).
            \]
            Comparing the coefficients of $x^3$ and $x$ gives
            \[
                p+q=-1,
                \qquad
                2q-4p=-8.
            \]
            Hence $p=1$ and $q=-2$. Comparing the coefficient of $x^2$ gives
            $a=-4$. Thus
            \[
                r_1,r_2
                =
                \frac{-1\pm i\sqrt7}{2}.
            \]
            Therefore,
            \[
                \boxed{
                    a=-4,
                    \qquad
                    \{r_1,r_2\}
                    =
                    \left\{
                        \frac{-1+i\sqrt7}{2},
                        \frac{-1-i\sqrt7}{2}
                    \right\}
                }.
            \]
        \end{expectedsubsolution}

        \begin{expectedsubproblem}
            Let $\alpha_1,\ldots,\alpha_5$ be the roots of
            \[
                x^5-5x+3=0.
            \]
            Determine
            \[
                \sum_{j=1}^{5}
                \left(\alpha_j^4+\alpha_j^5\right).
            \]
        \end{expectedsubproblem}

        \begin{expectedsubsolution}
            Let
            \[
                s_k=\sum_{j=1}^{5}\alpha_j^k.
            \]
            Newton's identities for
            \[
                x^5+0x^4+0x^3+0x^2-5x+3
            \]
            give
            \[
                s_1=s_2=s_3=0,
            \]
            \[
                s_4+4(-5)=0,
                \qquad
                s_5+5(3)=0.
            \]
            Hence $s_4=20$ and $s_5=-15$, so
            \[
                \boxed{s_4+s_5=5}.
            \]
        \end{expectedsubsolution}

        \begin{expectedsubproblem}
            Let $\alpha_1,\ldots,\alpha_4$ be the roots of
            \[
                x^4+x^3+x^2+x+1=0.
            \]
            Determine
            \[
                \sum_{j=1}^{4}\frac{1}{(1-\alpha_j)^2}.
            \]
        \end{expectedsubproblem}

        \begin{expectedsubsolution}
            Let
            \[
                p(x)=x^4+x^3+x^2+x+1
                =
                \prod_{j=1}^{4}(x-\alpha_j).
            \]
            Then
            \[
                \frac{p'(x)}{p(x)}
                =
                \sum_{j=1}^{4}\frac{1}{x-\alpha_j}.
            \]
            Differentiating gives
            \[
                \sum_{j=1}^{4}\frac{1}{(x-\alpha_j)^2}
                =
                \left(\frac{p'(x)}{p(x)}\right)^2
                -
                \frac{p''(x)}{p(x)}.
            \]
            Since
            \[
                p(1)=5,
                \qquad
                p'(1)=10,
                \qquad
                p''(1)=20,
            \]
            the required sum is
            \[
                \boxed{\left(\frac{10}{5}\right)^2-\frac{20}{5}=0}.
            \]
        \end{expectedsubsolution}

        \begin{expectedsubproblem}
            Let $\alpha_1,\alpha_2,\alpha_3$ be the roots of
            \[
                x^3-x-1=0.
            \]
            Determine
            \[
                \prod_{j=1}^{3}(1+\alpha_j^2).
            \]
        \end{expectedsubproblem}

        \begin{expectedsubsolution}
            Let
            \[
                p(x)=x^3-x-1
                =
                \prod_{j=1}^{3}(x-\alpha_j).
            \]
            For each root $\alpha_j$,
            \[
                (i-\alpha_j)(-i-\alpha_j)=1+\alpha_j^2.
            \]
            Therefore,
            \[
                \prod_{j=1}^{3}(1+\alpha_j^2)
                =
                p(i)p(-i).
            \]
            Since
            \[
                p(i)=-1-2i,
                \qquad
                p(-i)=-1+2i,
            \]
            we obtain
            \[
                \boxed{p(i)p(-i)=5}.
            \]
        \end{expectedsubsolution}

        \begin{expectedsubproblem}[1995 CSAT]
            Suppose that
            \[
                \frac{1}{(x-1)(x-2)\cdots(x-10)}
                =
                \sum_{k=1}^{10}\frac{a_k}{x-k}
            \]
            whenever the denominators are nonzero. Determine
            \[
                a_1+a_2+\cdots+a_{10}.
            \]
        \end{expectedsubproblem}

        \begin{expectedsubsolution}
            Multiply the identity by $x$ and let $x\to\infty$. The left-hand side
            tends to $0$, while
            \[
                \frac{x}{x-k}\to1
            \]
            for every $k$. Therefore
            \[
                \boxed{a_1+a_2+\cdots+a_{10}=0}.
            \]
        \end{expectedsubsolution}

        \begin{expectedsubproblem}
            Factor completely:
            \[
                (a-b)^3+(b-c)^3+(c-a)^3.
            \]
        \end{expectedsubproblem}

        \begin{expectedsubsolution}
            Set
            \[
                x=a-b,
                \qquad
                y=b-c,
                \qquad
                z=c-a.
            \]
            Then $x+y+z=0$. The identity
            \[
                x^3+y^3+z^3-3xyz
                =
                (x+y+z)
                (x^2+y^2+z^2-xy-yz-zx)
            \]
            therefore gives
            \[
                x^3+y^3+z^3=3xyz.
            \]
            Hence
            \[
                \boxed{3(a-b)(b-c)(c-a)}.
            \]
        \end{expectedsubsolution}

    \end{expectedproblemgroup}

    \begin{expectedproblemgroup}{Determinants}
        Compute each of the following determinants in exact form.

        \begin{expectedsubproblem}[2016 \KAIST{} Problem of the Week]
            Let $\mathbf{A}=(a_{ij})$ be the $n\times n$ matrix defined by
            \[
                a_{ij}
                =
                \begin{cases}
                    3, & i=j,\\
                    (-1)^{|i-j|}, & i\neq j.
                \end{cases}
            \]
            Determine $\det(\mathbf{A})$.
        \end{expectedsubproblem}

        \begin{expectedsubsolution}
            Let
            \[
                \mathbf{S}=\diag((-1)^1,(-1)^2,\ldots,(-1)^n).
            \]
            Then $\det(\mathbf{S})^2=1$, and
            \[
                \mathbf{S}\mathbf{A}\mathbf{S}
                =
                2\mathbf{I}_n+\mathbf{J}_n,
            \]
            where $\mathbf{J}_n$ is the all-ones matrix. The eigenvalues of
            $2\mathbf{I}_n+\mathbf{J}_n$ are $n+2$ once and $2$ with multiplicity
            $n-1$. Therefore,
            \[
                \boxed{\det(\mathbf{A})=2^{n-1}(n+2)}.
            \]
        \end{expectedsubsolution}

        \begin{expectedsubproblem}[2016 \KAIST{} Problem of the Week]
            Let $T$ be a tree on the vertex set $\{1,2,\ldots,n\}$, and let
            $d_{ij}$ be the distance between vertices $i$ and $j$. Define
            \[
                \mathbf{D}_T(x)=\left(x^{d_{ij}}\right)_{1\leq i,j\leq n}.
            \]
            Determine $\det(\mathbf{D}_T(x))$.
        \end{expectedsubproblem}

        \begin{expectedsubsolution}
            We use induction on $n$. The formula is immediate for $n=1$. For
            $n\geq2$, choose a leaf $n$ whose unique neighbor is $p$. For every
            $j\neq n$,
            \[
                d_{nj}=1+d_{pj}.
            \]
            Apply the row operation
            \[
                R_n\longleftarrow R_n-xR_p.
            \]
            The last row becomes zero except for its last entry $1-x^2$. Next apply
            \[
                C_n\longleftarrow C_n-xC_p.
            \]
            The last column also becomes zero except for the same last entry, while
            the remaining principal block is the matrix associated with the tree
            obtained by deleting the leaf. Thus
            \[
                \det(\mathbf{D}_T(x))
                =
                (1-x^2)\det(\mathbf{D}_{T-n}(x)).
            \]
            Induction gives
            \[
                \boxed{
                    \det(\mathbf{D}_T(x))=(1-x^2)^{n-1}
                }.
            \]
        \end{expectedsubsolution}

        \begin{expectedsubproblem}[2015 \KAIST{} Problem of the Week]
            For $n\in\N$ and $x\in\R$, let
            $\mathbf{T}_n(x)=(t_{ij})$ be the $n\times n$ tridiagonal matrix
            defined by
            \[
                t_{ij}
                =
                \begin{cases}
                    1, & i=j,\\
                    x, & |i-j|=1,\\
                    0, & |i-j|\geq2.
                \end{cases}
            \]
            Determine $\det(\mathbf{T}_n(x))$.
        \end{expectedsubproblem}

        \begin{expectedsubsolution}
            Put
            \[
                D_n(x)=\det(\mathbf{T}_n(x)),
                \qquad
                D_0(x)=1.
            \]
            Expansion along the final row gives
            \[
                D_n(x)=D_{n-1}(x)-x^2D_{n-2}(x),
                \qquad
                D_0(x)=D_1(x)=1.
            \]
            The polynomial
            \[
                \sum_{k=0}^{\lfloor n/2\rfloor}
                (-1)^k\binom{n-k}{k}x^{2k}
            \]
            has these initial values and satisfies the same recurrence by
            Pascal's identity. Therefore,
            \[
                \boxed{
                    \det(\mathbf{T}_n(x))
                    =
                    \sum_{k=0}^{\lfloor n/2\rfloor}
                    (-1)^k\binom{n-k}{k}x^{2k}
                }.
            \]
            Equivalently, if $x\neq\pm1/2$ and
            \[
                r_{\pm}=\frac{1\pm\sqrt{1-4x^2}}{2},
            \]
            then
            \[
                D_n(x)=\frac{r_+^{n+1}-r_-^{n+1}}{r_+-r_-};
            \]
            for $x=\pm1/2$, this becomes $D_n(x)=(n+1)/2^n$.
        \end{expectedsubsolution}

        \begin{expectedsubproblem}[2011 \KAIST{} Problem of the Week]
            Let $\mathbf{M}=(m_{ij})$ be the $n\times n$ matrix defined by
            \[
                m_{ij}=i(i+1)(i+2)\cdots(i+j-2),
            \]
            where the empty product $m_{i1}$ is $1$. Determine
            $\det(\mathbf{M})$.
        \end{expectedsubproblem}

        \begin{expectedsubsolution}
            For $j\geq1$, define the monic polynomial
            \[
                p_{j-1}(t)=t(t+1)\cdots(t+j-2),
            \]
            with $p_0(t)=1$. Then
            \[
                m_{ij}=p_{j-1}(i).
            \]
            Replacing the monomial basis
            $1,t,t^2,\ldots,t^{n-1}$ by the monic basis
            $p_0,p_1,\ldots,p_{n-1}$ uses a unit upper-triangular change-of-basis
            matrix. Hence $\det(\mathbf{M})$ equals the Vandermonde determinant at
            $1,2,\ldots,n$:
            \[
                \begin{aligned}
                    \det(\mathbf{M})
                    &=
                    \prod_{1\leq i<j\leq n}(j-i)\\
                    &=
                    \boxed{\prod_{k=1}^{n-1}k!}.
                \end{aligned}
            \]
        \end{expectedsubsolution}

        \begin{expectedsubproblem}[2012 \KAIST{} Problem of the Week]
            List all nonempty subsets of $\{1,2,\ldots,n\}$ as
            $S_1,S_2,\ldots,S_{2^n-1}$, and define the
            $(2^n-1)\times(2^n-1)$ matrix $\mathbf{A}=(a_{ij})$ by
            \[
                a_{ij}
                =
                \begin{cases}
                    1, & S_i\cap S_j\neq\emptyset,\\
                    0, & S_i\cap S_j=\emptyset.
                \end{cases}
            \]
            Determine $|\det(\mathbf{A})|$.
        \end{expectedsubproblem}

        \begin{expectedsubsolution}
            Index rows and columns by the nonempty subsets of $[n]$. Define
            \[
                c_{S,U}
                =
                \begin{cases}
                    1, & U\subseteq S,\\
                    0, & U\nsubseteq S,
                \end{cases}
            \]
            and let $\mathbf{C}=(c_{S,U})$. Also let $\mathbf{E}$ be diagonal with
            \[
                e_{U,U}=(-1)^{|U|+1}.
            \]
            Inclusion-exclusion gives
            \[
                (\mathbf{C}\mathbf{E}\mathbf{C}^{\mathsf{T}})_{S,T}
                =
                \sum_{\varnothing\neq U\subseteq S\cap T}(-1)^{|U|+1}
                =
                \begin{cases}
                    1, & S\cap T\neq\varnothing,\\
                    0, & S\cap T=\varnothing.
                \end{cases}
            \]
            Thus
            \[
                \mathbf{A}=\mathbf{C}\mathbf{E}\mathbf{C}^{\mathsf{T}}.
            \]
            Order the subsets by increasing cardinality. Then $\mathbf{C}$ is
            triangular with every diagonal entry equal to $1$, so
            $\det(\mathbf{C})=1$. Since every diagonal entry of $\mathbf{E}$ is
            $\pm1$,
            \[
                \boxed{|\det(\mathbf{A})|=1}.
            \]
        \end{expectedsubsolution}

    \end{expectedproblemgroup}

    \begin{expectedproblemgroup}{Eigenvalues and Matrix Powers}
        Determine the requested eigenvalues or matrix powers exactly. Use
        diagonalization, polynomial identities, or invariant subspaces whenever
        they make the computation shorter than direct multiplication.

        \begin{expectedsubproblem}[1997 VTRMC]
            Let
            \[
                \mathbf{A}
                =
                \begin{pmatrix}
                    1 & 3\\
                    1 & 1
                \end{pmatrix}.
            \]
            Determine the $(1,1)$-entry of $\mathbf{A}^{100}$.
        \end{expectedsubproblem}

        \begin{expectedsubsolution}
            The eigenvalues of $\mathbf{A}$ are
            \[
                \lambda_\pm=1\pm\sqrt3.
            \]
            The matrix is diagonalizable, and the two spectral projections have
            equal $(1,1)$-entries $1/2$. Hence
            \[
                (\mathbf{A}^{100})_{11}
                =
                \frac{\lambda_+^{100}+\lambda_-^{100}}{2}.
            \]
            Therefore,
            \[
                \boxed{
                    (\mathbf{A}^{100})_{11}
                    =
                    \frac{(1+\sqrt3)^{100}+(1-\sqrt3)^{100}}{2}
                }.
            \]
        \end{expectedsubsolution}

        \begin{expectedsubproblem}[Fibonacci Q-Matrix]
            Let
            \[
                \mathbf{Q}
                =
                \begin{pmatrix}
                    1 & 1\\
                    1 & 0
                \end{pmatrix}.
            \]
            Determine $\mathbf{Q}^n$ for $n\in\N$ in terms of the Fibonacci
            sequence.
        \end{expectedsubproblem}

        \begin{expectedsubsolution}
            Let $F_0=0$, $F_1=1$, and $F_{n+2}=F_{n+1}+F_n$. We claim that
            \[
                \mathbf{Q}^n
                =
                \begin{pmatrix}
                    F_{n+1} & F_n\\
                    F_n & F_{n-1}
                \end{pmatrix}
            \]
            for $n\geq1$. The formula is true for $n=1$, and multiplying the
            displayed matrix by $\mathbf{Q}$ advances every entry according to the
            Fibonacci recurrence. Thus
            \[
                \boxed{
                    \mathbf{Q}^n
                    =
                    \begin{pmatrix}
                        F_{n+1} & F_n\\
                        F_n & F_{n-1}
                    \end{pmatrix}
                }.
            \]
        \end{expectedsubsolution}


        \begin{expectedsubproblem}
            Let
            \[
                \mathbf{A}
                =
                \begin{pmatrix}
                    1 & 1 & 0\\
                    0 & 1 & 1\\
                    0 & 0 & 1
                \end{pmatrix}.
            \]
            Determine $\mathbf{A}^n$ for $n\in\Nzero$.
        \end{expectedsubproblem}

        \begin{expectedsubsolution}
            Write
            \[
                \mathbf{A}=\mathbf{I}_3+\mathbf{N},
                \qquad
                \mathbf{N}
                =
                \begin{pmatrix}
                    0 & 1 & 0\\
                    0 & 0 & 1\\
                    0 & 0 & 0
                \end{pmatrix}.
            \]
            Since $\mathbf{N}^3=\mathbf{0}$, the nilpotent-matrix recognition rule
            and the binomial theorem from Mathematical Concepts give
            \[
                \mathbf{A}^n
                =
                \mathbf{I}_3+n\mathbf{N}+\binom{n}{2}\mathbf{N}^2.
            \]
            Therefore,
            \[
                \boxed{
                    \mathbf{A}^n
                    =
                    \begin{pmatrix}
                        1 & n & \binom{n}{2}\\
                        0 & 1 & n\\
                        0 & 0 & 1
                    \end{pmatrix}
                }.
            \]
        \end{expectedsubsolution}

        \begin{expectedsubproblem}
            Determine the eigenvalues of
            \[
                \mathbf{A}
                =
                \begin{pmatrix}
                    1 & 2 & 3\\
                    1 & 2 & 3\\
                    1 & 2 & 3
                \end{pmatrix}.
            \]
        \end{expectedsubproblem}

        \begin{expectedsubsolution}
            Write
            \[
                \mathbf{A}
                =
                \begin{pmatrix}
                    1\\1\\1
                \end{pmatrix}
                \begin{pmatrix}
                    1&2&3
                \end{pmatrix}.
            \]
            This is a rank-one outer product. Its only possible nonzero
            eigenvalue is
            \[
                \begin{pmatrix}
                    1&2&3
                \end{pmatrix}
                \begin{pmatrix}
                    1\\1\\1
                \end{pmatrix}
                =
                6.
            \]
            Since the matrix has rank $1$, the remaining two eigenvalues are
            $0$. Therefore,
            \[
                \boxed{6,\ 0,\ 0}.
            \]
        \end{expectedsubsolution}

        \begin{expectedsubproblem}
            Determine the eigenvalues of
            \[
                \mathbf{A}
                =
                \begin{pmatrix}
                    1 & -2 & 3 & 0\\
                    -2 & 3 & 0 & 1\\
                    3 & 0 & 1 & -2\\
                    0 & 1 & -2 & 3
                \end{pmatrix}.
            \]
        \end{expectedsubproblem}

        \begin{expectedsubsolution}
            Reorder the coordinates as $(1,3,2,4)$. The resulting similar
            matrix has the block form
            \[
                \begin{pmatrix}
                    \mathbf{B} & -2\mathbf{I}_2\\
                    -2\mathbf{I}_2 & \mathbf{C}
                \end{pmatrix},
                \qquad
                \mathbf{B}
                =
                \begin{pmatrix}
                    1&3\\3&1
                \end{pmatrix},
                \qquad
                \mathbf{C}
                =
                \begin{pmatrix}
                    3&1\\1&3
                \end{pmatrix}.
            \]
            The vectors $(1,1)$ and $(1,-1)$ simultaneously diagonalize
            $\mathbf{B}$ and $\mathbf{C}$. On the corresponding two-dimensional
            invariant subspaces, the matrix reduces respectively to
            \[
                \begin{pmatrix}
                    4&-2\\-2&4
                \end{pmatrix}
                \quad\text{and}\quad
                \begin{pmatrix}
                    -2&-2\\-2&2
                \end{pmatrix}.
            \]
            Their eigenvalues are $6,2$ and $\pm2\sqrt2$. Hence
            \[
                \boxed{6,\ 2,\ 2\sqrt2,\ -2\sqrt2}.
            \]
        \end{expectedsubsolution}

        \begin{expectedsubproblem}
            Let $\mathbf{A}$ be the $2022\times2022$ matrix whose diagonal
            entries are $0$ and whose off-diagonal entries are $1$. Determine
            all eigenvalues, including multiplicities.
        \end{expectedsubproblem}

        \begin{expectedsubsolution}
            We have
            \[
                \mathbf{A}=\mathbf{J}_{2022}-\mathbf{I}_{2022}.
            \]
            The all-ones vector is an eigenvector of $\mathbf{J}_{2022}$ with
            eigenvalue $2022$, while every vector whose coordinates sum to zero
            has eigenvalue $0$. Subtracting the identity shifts all eigenvalues
            by $-1$. Therefore,
            \[
                \boxed{
                    2021\text{ with multiplicity }1,
                    \qquad
                    -1\text{ with multiplicity }2021
                }.
            \]
        \end{expectedsubsolution}


        \begin{expectedsubproblem}
            Determine all eigenvalues of
            \[
                \mathbf{P}
                =
                \begin{pmatrix}
                    0&1&0&0\\
                    0&0&1&0\\
                    0&0&0&1\\
                    1&0&0&0
                \end{pmatrix}.
            \]
        \end{expectedsubproblem}

        \begin{expectedsubsolution}
            The matrix cyclically shifts the coordinates and satisfies
            \[
                \mathbf{P}^4=\mathbf{I}_4.
            \]
            Hence every eigenvalue $\lambda$ satisfies $\lambda^4=1$. The
            characteristic polynomial is
            \[
                \lambda^4-1,
            \]
            whose roots are
            \[
                1,\quad -1,\quad i,\quad -i.
            \]
            Therefore the eigenvalues are
            \[
                \boxed{1,\ -1,\ i,\ -i}.
            \]
        \end{expectedsubsolution}

        \begin{expectedsubproblem}
            Let $\mathbf{A}=(a_{ij})$ be the $n\times n$ matrix
            \[
                a_{ij}=ij,
                \qquad
                1\leq i,j\leq n.
            \]
            Determine the sum of all eigenvalues of $\mathbf{A}$, counted with
            algebraic multiplicity.
        \end{expectedsubproblem}

        \begin{expectedsubsolution}
            Mathematical Concepts gives that the sum of the eigenvalues,
            counted with algebraic multiplicity, equals the trace. Therefore
            \[
                \sum_{k=1}^{n}\lambda_k
                =
                \tr(\mathbf{A})
                =
                \sum_{i=1}^{n}i^2
                =
                \boxed{\frac{n(n+1)(2n+1)}{6}}.
            \]
            Equivalently, $\mathbf{A}=\mathbf{u}\mathbf{u}^{\mathsf T}$ for
            $\mathbf{u}=(1,2,\ldots,n)^{\mathsf T}$, so $\mathbf{A}$ has rank
            $1$ and its unique nonzero eigenvalue is $\|\mathbf{u}\|^2$.
        \end{expectedsubsolution}

        \begin{expectedsubproblem}
            Let
            \[
                \mathbf{A}
                =
                \begin{pmatrix}
                    -11 & 4\\
                    -30 & 11
                \end{pmatrix}.
            \]
            Determine $\mathbf{A}^{2016}$.
        \end{expectedsubproblem}

        \begin{expectedsubsolution}
            A direct multiplication gives
            \[
                \mathbf{A}^2
                =
                \begin{pmatrix}
                    1&0\\
                    0&1
                \end{pmatrix}
                =
                \mathbf{I}_2.
            \]
            Hence
            \[
                \mathbf{A}^{2016}
                =
                (\mathbf{A}^2)^{1008}
                =
                \boxed{\mathbf{I}_2}.
            \]
        \end{expectedsubsolution}

    \end{expectedproblemgroup}

